\chapter{从一个问题出发：调控怎样改变表达的时间行为}
一个以恒定速率产生、以一阶速率降解的蛋白质，总会逐渐回到同一个浓度。基因调控的作用，是让产生率依赖其他分子，甚至依赖自身：同样的外部条件下，系统于是可能保留不同表达状态，也可能持续振荡，或者只对输入变化作出短暂响应。要理解这些现象，需要依次回答三个问题：产生率从什么反应机制来？反馈如何改变稳态和扰动的演化？这些局部行为在随机性与网络背景中还能否成立？

每个模型依次讲解回路、方程与适用条件，比较新增环节怎样带来状态选择、振荡或适应，再回到分子计数和真实网络。

网络节点常表示基因表达产物（例如蛋白质）的有效浓度。启动子状态、RNA 聚合酶、mRNA、配体和代谢物等微观变量可能在简化时被省略。本文用 $x,y,z$ 表示按参考浓度换算后的数值，例如 $x=1$ 表示达到参考浓度。激活边表示：调控因子增加时，目标的产生率上升；抑制边表示产生率下降。用公式说，就是 $\partial F_i/\partial x_j$ 分别为正或负。边的符号说明作用方向，具体结合方式还要另外建模。

\begin{figure}[H]\centering
\begin{tikzpicture}[node distance=22mm]
\node[gene] (x) {$X$}; \node[gene,right=of x] (y) {$Y$};
\draw[act] (x)--node[above]{激活：$+$}(y);
\node[gene,right=35mm of y] (u) {$X$}; \node[gene,right=of u] (v) {$Y$};
\draw[rep] (u)--node[above]{抑制：$-$}(v);
\end{tikzpicture}
\caption{调控边的记号。箭头表示产生率随调控变量增加，钝头表示产生率下降；边的符号不能单独指定结合机制。}\label{fig:signs}
\end{figure}

标准连续模型为
\begin{equation}\dot x_i=F_i(\bm{x};\bm{u})-\gamma_i x_i,\qquad F_i\geq0,\quad\gamma_i>0.\end{equation}
其中 $F_i$ 为有效产生率，$\bm u$ 为外部输入，$\gamma_i$ 汇总一阶降解及相应的稀释效应。讨论稳态时固定输入；讨论响应时再指定输入怎样变化。同一组初始浓度能否决定唯一一条轨迹，是后面所有相图分析的前提。一个常用的保证条件叫\emph{局部 Lipschitz 连续}：在每个所考察状态附近，变化率 $\bm v(\bm x)$ 满足
\[
\|\bm v(\bm x)-\bm v(\bm y)\|\leq L\|\bm x-\bm y\|,
\]
其中 $L$ 是局部范围内的有限常数，$\|\cdot\|$ 表示状态之间的距离。直观地说，浓度只改动一点，变化率也只能改动有限的一点，不能在很小的浓度差上出现无限陡的变化。这个条件比“曲线不断开”更强，却不要求每个位置都有斜率；它保证同一起点不会分叉出两种未来。例如 $v(x)=|x|$ 满足 $|v(x)-v(y)|\leq|x-y|$，可取 $L=1$；虽然零点处有尖角、没有导数，它仍然是 Lipschitz 连续的。“局部”表示不同浓度范围可以用不同的 $L$，并不是要求整个浓度空间共用一个限制。

使用导数分析小扰动时，本文通常进一步假设产生率为 $C^1$。这个记号表示\emph{函数可求一阶导数，且一阶导数连续}：不仅曲线不断开，斜率本身也连续变化。$C^1$ 函数在足够小的范围内满足上述 Lipschitz 条件，因此本文的 Hill 等常见模型可以同时支持唯一轨迹和斜率分析。这些是假设在所分析的浓度范围内成立；若模型有开关式跳变或边界处无限陡的斜率，就要重新检查。

\emph{稳态}指产生与损失平衡、浓度不再变化。\emph{稳定}指小偏差下轨迹始终留在附近；若偏差逐渐消失，则称为\emph{渐近稳定}。局部结论只涉及稳态附近，不代表所有初始浓度都回到它。
